% Source: Weinberg GR, physical PDF pp. 208--211
%         (printed pp. 185--188).
% Coverage: Section 8.4, General Equations of Motion; equations
%           (8.4.1)--(8.4.30).
% Figures/tables/footnotes: none.
% Status: source-reviewed and compile-clean.
% Uncertainties: Weinberg's trajectory parameter p is modernized to the
%                conventional affine parameter lambda throughout.

\subsection{General Equations of Motion}\label{sec:8.4}

We now consider the motion of a freely falling material particle or photon
in a static isotropic gravitational field. First let us consider the most
general such metric in the standard form derived in Section~\ref{sec:8.1},
that is,
\begin{equation}
  \dd s^2
  =
  -B(r)\,\dd t^2+A(r)\,\dd r^2
  +r^2\dd\theta^2+r^2\sin^2\theta\,\dd\varphi^2.
  \tag{8.4.1}\label{eq:8.4.1}
\end{equation}
The equations of free fall are
\begin{equation}
  \frac{\dd^2x^\mu}{\dd\lambda^2}
  +\Gamma^\mu{}_{\nu\rho}
  \frac{\dd x^\nu}{\dd\lambda}\frac{\dd x^\rho}{\dd\lambda}
  =0,
  \tag{8.4.2}\label{eq:8.4.2}
\end{equation}
where \(\lambda\) is an affine parameter describing the trajectory. In
general \(\dd\tau\) is proportional to \(\dd\lambda\), so for a material
particle we could normalize \(\lambda\) so that \(\lambda=\tau\). However,
for a photon the proportionality constant \(\dd\tau/\dd\lambda\) vanishes, and
since we wish to treat photons as well as massive particles, we shall find
it convenient to reserve the right to fix the normalization of \(\lambda\)
independently from that of \(\tau\).

Using the nonvanishing components of the affine connection given by
Eq.~\eqref{eq:8.1.11}, we find from Eq.~\eqref{eq:8.4.2} that
\begin{equation}
  \begin{split}
  0={}&
  \frac{\dd^2r}{\dd\lambda^2}
  +\frac{A'(r)}{2A(r)}
    \left(\frac{\dd r}{\dd\lambda}\right)^2
  -\frac{r}{A(r)}
    \left(\frac{\dd\theta}{\dd\lambda}\right)^2\\
  &-\frac{r\sin^2\theta}{A(r)}
    \left(\frac{\dd\varphi}{\dd\lambda}\right)^2
  +\frac{B'(r)}{2A(r)}
    \left(\frac{\dd t}{\dd\lambda}\right)^2,
  \end{split}
  \tag{8.4.3}\label{eq:8.4.3}
\end{equation}
\begin{equation}
  0
  =
  \frac{\dd^2\theta}{\dd\lambda^2}
  +\frac{2}{r}\frac{\dd\theta}{\dd\lambda}\frac{\dd r}{\dd\lambda}
  -\sin\theta\cos\theta
  \left(\frac{\dd\varphi}{\dd\lambda}\right)^2,
  \tag{8.4.4}\label{eq:8.4.4}
\end{equation}
\begin{equation}
  0
  =
  \frac{\dd^2\varphi}{\dd\lambda^2}
  +\frac{2}{r}\frac{\dd\varphi}{\dd\lambda}\frac{\dd r}{\dd\lambda}
  +2\cot\theta
  \frac{\dd\varphi}{\dd\lambda}\frac{\dd\theta}{\dd\lambda},
  \tag{8.4.5}\label{eq:8.4.5}
\end{equation}
\begin{equation}
  0
  =
  \frac{\dd^2t}{\dd\lambda^2}
  +\frac{B'(r)}{B(r)}
  \frac{\dd t}{\dd\lambda}\frac{\dd r}{\dd\lambda}.
  \tag{8.4.6}\label{eq:8.4.6}
\end{equation}
(A prime denotes \(d/\dd r\).) We solve these equations by looking for
constants of the motion.

Since the field is isotropic, we may consider the orbit of our particle to
be confined to the equatorial plane, that is,
\begin{equation}
  \theta=\frac{\pi}{2}.
  \tag{8.4.7}\label{eq:8.4.7}
\end{equation}
Then Eq.~\eqref{eq:8.4.4} is immediately satisfied and we can forget about
\(\theta\) as a dynamical variable. Dividing Eqs.~\eqref{eq:8.4.5} and
\eqref{eq:8.4.6} by \(\dd\varphi/\dd\lambda\) and \(\dd t/\dd\lambda\),
respectively, we next find
\begin{equation}
  \frac{\dd}{\dd\lambda}
  \left[
    \ln\frac{\dd\varphi}{\dd\lambda}+\ln r^2
  \right]=0,
  \tag{8.4.8}\label{eq:8.4.8}
\end{equation}
\begin{equation}
  \frac{\dd}{\dd\lambda}
  \left[
    \ln\frac{\dd t}{\dd\lambda}+\ln B
  \right]=0.
  \tag{8.4.9}\label{eq:8.4.9}
\end{equation}
This yields two constants of the motion. One of them will be absorbed
immediately into the definition of \(\lambda\); we choose to normalize
\(\lambda\) so that the solution of Eq.~\eqref{eq:8.4.9} is
\begin{equation}
  \frac{\dd t}{\dd\lambda}=\frac{1}{B(r)}.
  \tag{8.4.10}\label{eq:8.4.10}
\end{equation}
Since \(B(r)\) is close to unity, \(\lambda\) is nearly equal to the
coordinate time \(t\). The other constant is obtained from
Eq.~\eqref{eq:8.4.8}, and plays the role of an angular momentum per unit
mass:
\begin{equation}
  r^2\frac{\dd\varphi}{\dd\lambda}=J
  \qquad\text{(constant)}.
  \tag{8.4.11}\label{eq:8.4.11}
\end{equation}
Inserting Eqs.~\eqref{eq:8.4.7}, \eqref{eq:8.4.10}, and
\eqref{eq:8.4.11} in Eq.~\eqref{eq:8.4.3} gives the remaining equation
of motion as
\begin{equation}
  0
  =
  \frac{\dd^2r}{\dd\lambda^2}
  +\frac{A'(r)}{2A(r)}
    \left(\frac{\dd r}{\dd\lambda}\right)^2
  -\frac{J^2}{r^3A(r)}
  +\frac{B'(r)}{2A(r)B^2(r)}.
  \tag{8.4.12}\label{eq:8.4.12}
\end{equation}
By multiplying this equation with \(2A(r)\,\dd r/\dd\lambda\), we may write it
as
\[
  \frac{\dd}{\dd\lambda}
  \left[
    A(r)\left(\frac{\dd r}{\dd\lambda}\right)^2
    +\frac{J^2}{r^2}-\frac{1}{B(r)}
  \right]=0,
\]
and our last constant of the motion is therefore
\begin{equation}
  A(r)\left(\frac{\dd r}{\dd\lambda}\right)^2
  +\frac{J^2}{r^2}-\frac{1}{B(r)}
  =-E
  \qquad\text{(constant)}.
  \tag{8.4.13}\label{eq:8.4.13}
\end{equation}
The proper time \(\tau\) may now be determined from
Eqs.~\eqref{eq:8.4.1}, \eqref{eq:8.4.7}, \eqref{eq:8.4.10},
\eqref{eq:8.4.11}, and \eqref{eq:8.4.13}; we find that
\begin{equation}
  \dd\tau^2=E\,\dd\lambda^2,
  \tag{8.4.14}\label{eq:8.4.14}
\end{equation}
in accordance with our earlier remark that Eq.~\eqref{eq:8.4.2} forces
\(\dd\tau/\dd\lambda\) to be constant. We see that \(E\) must take the values
\begin{equation}
  E>0\qquad\text{for material particles},
  \tag{8.4.15}\label{eq:8.4.15}
\end{equation}
\begin{equation}
  E=0\qquad\text{for photons}.
  \tag{8.4.16}\label{eq:8.4.16}
\end{equation}
Also, \(A(r)\) is in practice always positive, so
Eq.~\eqref{eq:8.4.13} tells us that our particle can reach radius \(r\)
only if
\begin{equation}
  \frac{J^2}{r^2}+E\leq\frac{1}{B(r)}.
  \tag{8.4.17}\label{eq:8.4.17}
\end{equation}

The parameter \(\lambda\) may be eliminated everywhere by using
Eq.~\eqref{eq:8.4.10} in Eqs.~\eqref{eq:8.4.11},
\eqref{eq:8.4.13}, and \eqref{eq:8.4.14}; we then have
\begin{equation}
  r^2\frac{\dd\varphi}{\dd t}=JB(r),
  \tag{8.4.18}\label{eq:8.4.18}
\end{equation}
\begin{equation}
  \frac{A(r)}{B^2(r)}
  \left(\frac{\dd r}{\dd t}\right)^2
  +\frac{J^2}{r^2}
  -\frac{1}{B(r)}
  =-E,
  \tag{8.4.19}\label{eq:8.4.19}
\end{equation}
\begin{equation}
  \dd\tau^2=EB^2(r)\,\dd t^2.
  \tag{8.4.20}\label{eq:8.4.20}
\end{equation}
For a slowly moving particle in a weak field \(J^2/r^2\),
\((\dd r/\dd t)^2\), \(A-1\), and \(B-1\simeq2\phi\) will all be small, and
to first order in these quantities the above equations of motion become
\[
  r^2\frac{\dd\varphi}{\dd t}\simeq J,
\]
\[
  \frac12\left(\frac{\dd r}{\dd t}\right)^2
  +\frac{J^2}{2r^2}+\phi
  \simeq\frac{1-E}{2}.
\]
These are the same equations as would hold in Newton's theory, with
\((1-E)/2\) playing the role of an energy per unit mass.

To see how the exact equations of motion work in a simple case, consider a
particle in a circular orbit at radius \(R\). Since \(\dd r/\dd t\) vanishes,
Eq.~\eqref{eq:8.4.19} gives
\begin{equation}
  \frac{J^2}{R^2}-\frac{1}{B(R)}+E=0.
  \tag{8.4.21}\label{eq:8.4.21}
\end{equation}
Also, for equilibrium at this radius, the derivative at \(R\) of the
left-hand side must also vanish, so
\begin{equation}
  -\frac{2J^2}{R^3}+\frac{B'(R)}{B^2(R)}=0.
  \tag{8.4.22}\label{eq:8.4.22}
\end{equation}
(If we regard a circle as the limit of an ellipse with perihelia
\(R-\delta\) and aphelia \(R+\delta\), then Eq.~\eqref{eq:8.4.19}
shows that \(J^2/r^2-1/B(r)+E\) must vanish at \(r=R\pm\delta\), and
this gives Eqs.~\eqref{eq:8.4.21} and \eqref{eq:8.4.22} in the limit
\(\delta\to0\).) From Eqs.~\eqref{eq:8.4.21} and \eqref{eq:8.4.22}
we find
\begin{equation}
  E
  =
  \frac{1}{B(R)}
  \left(1-\frac{RB'(R)}{2B(R)}\right),
  \tag{8.4.23}\label{eq:8.4.23}
\end{equation}
\begin{equation}
  J^2=\frac{B'(R)R^3}{2B^2(R)}.
  \tag{8.4.24}\label{eq:8.4.24}
\end{equation}
Using Eq.~\eqref{eq:8.4.24} in Eq.~\eqref{eq:8.4.18} gives the rate of
revolution as
\begin{equation}
  \frac{\dd\varphi}{\dd t}
  =
  \left(\frac{B'(R)}{2R}\right)^{1/2},
  \tag{8.4.25}\label{eq:8.4.25}
\end{equation}
whereas Eqs.~\eqref{eq:8.4.23} and \eqref{eq:8.4.20} give the proper
time as
\begin{equation}
  \frac{\dd\tau}{\dd t}
  =
  \sqrt{B(R)-\frac12RB'(R)}.
  \tag{8.4.26}\label{eq:8.4.26}
\end{equation}
By using the Robertson expansion \eqref{eq:8.3.8}, we find
\begin{equation}
  \frac{\dd\varphi}{\dd t}
  =
  \left(\frac{MG}{R^3}\right)^{1/2}
  \left[1-(\beta-\gamma)\frac{MG}{R}+\cdots\right],
  \tag{8.4.27}\label{eq:8.4.27}
\end{equation}
\begin{equation}
  \frac{\dd\tau}{\dd t}
  =
  \left[1-\frac{3MG}{R}+\cdots\right].
  \tag{8.4.28}\label{eq:8.4.28}
\end{equation}

In most applications of general relativity we are more interested in the
shape of orbits, that is, in \(r\) as a function of \(\varphi\), than in
their time history. The orbit shape can be obtained directly by eliminating
\(\dd\lambda\) from Eqs.~\eqref{eq:8.4.11} and \eqref{eq:8.4.13}; this
gives
\begin{equation}
  \frac{A(r)}{r^4}
  \left(\frac{\dd r}{\dd\varphi}\right)^2
  +\frac{1}{r^2}
  -\frac{1}{J^2B(r)}
  =-\frac{E}{J^2}.
  \tag{8.4.29}\label{eq:8.4.29}
\end{equation}
The solution may then be determined by a quadrature:
\begin{equation}
  \varphi
  =
  \pm\int
  \frac{A^{1/2}(r)\,\dd r}
  {r^2\left[J^{-2}B^{-1}(r)-EJ^{-2}-r^{-2}\right]^{1/2}}.
  \tag{8.4.30}\label{eq:8.4.30}
\end{equation}
